\chapter{PROPUESTA Y METODOLOGÍA}\label{cap:propuesta}
% Bloque (c), parte formal y diseno experimental. Ver redaccion/MAPA.md.
% Pautas de la plantilla UTEC movidas a redaccion/RUBRICA.md (P-MM).
% Fuentes: tesis/main.tex (Objectives, Expected Results, Pre-registration, Datasets),
% experiments/objetivo2/preinscripcion_muestreador.md, docs/02-datos.md,
% docs/01-decisiones.md. Fuente de cada cifra: redaccion/rondas/capitulo3-r00-respuesta.md
% capitulo3-r01-respuesta.md a capitulo3-r05-respuesta.md.

Este capítulo describe la contribución y el diseño con que se evalúa. La contribución no es un generador único, sino una cadena de dos componentes que se diseñan y se evalúan por separado. El \textbf{muestreador de colocación} decide dónde y con qué pose tridimensional se ubica un tornillo transilíaco-transsacro (transiliosacro). El \textbf{sintetizador} genera la apariencia del implante y de su artefacto metálico local en la tomografía computarizada (TC). Antes de ambos se ubica una compuerta que decide en qué representación puede operar el sintetizador. El capítulo sigue ese orden: visión general, datos, compuerta de representación, muestreador, sintetizador, protocolo de evaluación y amenazas a la validez.

\section{Visión general}\label{sec:vision}

La propuesta se organiza en cuatro objetivos específicos \GAPDEC{alinear la pregunta de investigación y los objetivos de la introducción con los cuatro objetivos de este capítulo, que hoy no coinciden}. El Objetivo~1 evalúa si una representación de la TC conserva las unidades Hounsfield (HU) del hueso con una tolerancia fijada de antemano. El Objetivo~2 diseña el muestreador de colocación y lo evalúa con la admisibilidad quirúrgica de la colocación (SAP, \emph{Surgical Admissibility of Placement}), que es la única métrica que introduce este trabajo. El Objetivo~3 diseña el sintetizador condicionado. El Objetivo~4 formaliza SAP y adopta, para la apariencia, un protocolo publicado con sus métricas y sus nombres originales, sin renombrarlas como contribución.

La cadena recibe un volumen de TC pélvica sin osteosíntesis y devuelve el mismo volumen con un tornillo sintético y su artefacto. El muestreador mide sobre cada volumen un corredor óseo, perturba su eje y entrega una pose; esa pose, junto con un modelo paramétrico del tornillo, define la \textbf{máscara del implante $M$}. Alrededor de $M$ se añade una \textbf{banda de generación extendida $B_{\delta}$} (\emph{generation band}), y ambas forman la \textbf{región de generación} $G = M \cup B_{\delta}$. El sintetizador genera valores de TC solo dentro de $G$ y copia el resto del volumen de origen. La Figura~\ref{fig:pipeline} resume ese flujo.

\begin{figure}[htbp]
\centering
\small
\fbox{\parbox{0.8\linewidth}{\centering\itshape Compuerta de representación (Objetivo~1), previa a la cadena: su resultado negativo descartó la ruta latente\\ y fijó el dominio de imagen para el sintetizador}}\\[8pt]
\fbox{\parbox{0.8\linewidth}{\centering TC receptora sin osteosíntesis (subconjunto local de CTPelvic1K)}}\\[2pt]
$\downarrow$\\[2pt]
\fbox{\parbox{0.8\linewidth}{\centering Máscaras anatómicas (TotalSegmentator) y medición del corredor óseo\\ en el marco de referencia de Kaiser et al.}}\\[2pt]
$\downarrow$\\[2pt]
\fbox{\parbox{0.8\linewidth}{\centering \textbf{Muestreador de colocación} (Objetivo~2): poses perturbadas alrededor del eje del corredor,\\ evaluadas con SAP frente a las dos distribuciones de la referencia clínica}}\\[2pt]
$\downarrow$\\[2pt]
\fbox{\parbox{0.8\linewidth}{\centering Tornillo paramétrico en la pose muestreada: máscara $M$ y región $G = M \cup B_{\delta}$}}\\[2pt]
$\downarrow$\\[2pt]
\fbox{\parbox{0.8\linewidth}{\centering \textbf{Sintetizador} (Objetivo~3): \emph{inpainting} por difusión en el dominio de imagen dentro de $G$;\\ fuera de $G$ se copia el volumen de origen}}\\[2pt]
$\downarrow$\\[2pt]
\fbox{\parbox{0.8\linewidth}{\centering Evaluación de apariencia con las métricas de Peters et al., frente a la inserción por copia y pegado\\ y frente al protocolo físico}}
\caption{Cadena propuesta. Las cajas en negrita son los dos componentes que se diseñan y se evalúan por separado; la caja superior es la compuerta que fijó el dominio del sintetizador.}
\label{fig:pipeline}
\end{figure}

Lo que se evalúa es la coherencia física y quirúrgica de lo sintetizado, no su utilidad para entrenar redes de segmentación. Quedan fuera de alcance las ablaciones restricción por restricción y la evaluación de segmentación posterior (Dice, HD95). Esta última se excluye por dos condiciones de los datos. La primera es la anotación: de los 75 volúmenes del subconjunto CLINIC-metal de CTPelvic1K, la publicación anota 14 \cite{liu2021ctpelvic1k}, y no está verificado que esas máscaras estén disponibles localmente. La segunda es la cobertura: localmente se dispone de 178 de los 1\,184 volúmenes publicados (Sección~\ref{sec:datos}). Un resultado de segmentación descansaría, por tanto, en una cohorte pequeña y parcialmente anotada.

El estado de ejecución de cada objetivo condiciona lo que este capítulo puede describir como hecho. La compuerta del Objetivo~1 se ejecutó y dio un veredicto negativo; el muestreador y SAP del Objetivo~2 se ejecutaron sobre las cohortes primaria y de sensibilidad. El sintetizador del Objetivo~3 tiene formulación y conjunto de entrenamiento definidos \GAPDATO{la cadena completa, de la pose sobre el corredor medido y el rasterizado de $M$ a la generación de la apariencia sobre una pelvis sin osteosíntesis, se ejecutó sobre dos pacientes de validación; no existe aún resultado del Objetivo~3 porque el modelo final no está elegido (Sección~\ref{sec:sintetizador})}. El Objetivo~3 se rediseñó después de conocer el veredicto del Objetivo~1; la regla de la compuerta se fijó antes de correr la prueba que decide y no se modificó.

\section{Datos y auditoría de la cohorte}\label{sec:datos}

La fuente de anatomía y de observaciones reales de artefacto es CTPelvic1K, que declara 1\,184 volúmenes de TC pélvica \cite{liu2021ctpelvic1k}. Localmente se dispone de 178 de ellos, repartidos en dos carpetas: \texttt{dataset6}, con 103 volúmenes, y \texttt{dataset7}, con 75. La correspondencia de esas carpetas con los subconjuntos CLINIC y CLINIC-metal no está declarada en los encabezados de los archivos: se infiere por coincidencia de conteo, espaciado y dimensiones con la tabla de la publicación. Los recuentos de este trabajo describen el inventario local, no la colección completa.

El artículo del conjunto no reporta los parámetros de adquisición ni de reconstrucción, incluida una posible reducción de artefactos metálicos (MAR, \emph{metal artifact reduction}) del fabricante o una reconstrucción monoenergética. Selles et al.~\cite{selles2024marreview} describen que ambas cambian la apariencia del artefacto. En consecuencia, las observaciones reales de artefacto, cuya reconstrucción es desconocida, se usan para comparar la apariencia y no como verdad física.

El nombre de un subconjunto no establece el tipo de implante ni la disponibilidad de anotación, de modo que la cohorte se construyó por auditoría. Un cribado por HU prioriza la revisión: se marca todo volumen con al menos un vóxel sobre 2500~HU, sin filtro de tamaño. Ese umbral es una heurística de cribado, no un criterio validado de identificación, y no forma parte de las métricas de Peters et al. La autora revisó en vistas tridimensionales los 178 volúmenes y anotó tipo y cantidad de material; la elegibilidad para entrenamiento exige además revisión multiplanar de cortes \GAPDATO{constancia de que la revisión multiplanar completa de cortes está hecha; el registro de datos solo acredita la revisión tridimensional}.

La unidad de independencia es el paciente, no el volumen. Los duplicados exactos se detectan por la huella criptográfica (SHA256) del volumen completo y los parciales por la huella de cada corte, porque la primera no ve subconjuntos ni campos de visión solapados del mismo estudio. Esas dos huellas, más una observación visual que unió dos adquisiciones del mismo implante, agrupan los 178 volúmenes en 168 pacientes. Los 168 pacientes se reparten en tres grupos: 65 con material ortopédico (grupo~1), 37 con objetos no ortopédicos (grupo~2) y 66 sin objeto metálico (grupo~3). Los objetos no ortopédicos son, por ejemplo, material extracorpóreo o dispositivos intrauterinos. La partición trabaja sobre 179 unidades de volumen, porque dos campos de visión solapados del mismo estudio se unen sin interpolar en una unidad adicional. De esas unidades, 10 quedan fuera de uso y una se reserva como par de reproducibilidad del mismo paciente.

La \textbf{regla de aislamiento estricto} prohíbe que un paciente contribuya a la vez al entrenamiento y a la evaluación. En el Objetivo~1, la partición se hizo por paciente y estratificada por subconjunto y por presencia de metal. Separa tres conjuntos: uno de entrenamiento, uno de validación con 8 casos y uno de prueba con 34 pacientes. El Objetivo~3 reutiliza esa partición: sus pares de entrenamiento se forman solo con volúmenes de pacientes de entrenamiento con metal. El Objetivo~2 no entrena ningún modelo, por lo que su cohorte no se divide; se construye con pacientes sin osteosíntesis (grupos~2 y~3), porque el corredor debe medirse sobre hueso sin implante.

La cohorte del Objetivo~2 resulta de un embudo de dos pasos. De 103 pacientes sin osteosíntesis, 91 tenían localizado el nivel S1 y 72 superaron el control de nivel; esos 72 forman la cohorte primaria. El control de nivel excluye un caso en dos situaciones. La primera es que el nivel S1 de las máscaras anatómicas (Sección~\ref{sec:corredor}) no coincida con el que da la localización automática de referencias anatómicas (Sección~\ref{sec:marco-metal}). Esa comparación enfrenta el techo de la máscara de S1, etiqueta de vértebra completa que incluye el arco posterior, con el punto del platillo superior que da la localización. La segunda es que la máscara de S1 toque el borde del campo de visión o esté vacía. La discordancia de nivel excluyó a 11 de 34 pacientes con objetos no ortopédicos y a 7 de 57 pacientes sin objeto, y un paciente más se excluyó porque S1 tocaba el borde del campo de visión. La cohorte de sensibilidad son los 49 pacientes sin objeto que superan el mismo control.

El criterio de selección fue la ausencia de osteosíntesis y no la ausencia de fractura, de modo que la fractura se cribó después. Un cribado ciego de fractura sobre 30 casos, 15 de corredor estrecho y 15 de control, halló fractura confirmada en 20 de los 30 (67\,\%), más un caso dudoso que se reporta aparte y no se cuenta como fractura. La proporción es idéntica en los dos grupos: 10 de 15 corredores estrechos y 10 de 15 controles. El cribado lo hizo un médico recién egresado y sin especialidad, sin conocer a qué grupo pertenecía cada caso; es un cribado válido y no una lectura de especialista, y así se usa en este trabajo. Juzgó 18 casos sobre láminas fijas y 12 sobre el volumen completo, que abrió en un visor cuando la lámina no bastaba. Ese paso al volumen lo motivó una sospecha, de modo que la detección pudo ser más sensible donde ya se sospechaba algo, pero se repartió de forma pareja entre los grupos: 7 corredores estrechos y 5 controles. La cohorte receptora no es, por tanto, una colección de pelvis sanas. Ninguna cifra del corredor depende de ello, porque el corredor se mide sobre cada volumen tal como está, pero sí cambia cómo se compara esa anatomía con la de la referencia clínica (Sección~\ref{sec:amenazas}).

Los 65 pacientes con material ortopédico cumplen cuatro funciones y ninguna de ellas es la de aportar geometría de implante. Sirven de contraste descriptivo para las colocaciones, de comparación de apariencia para la síntesis, de fuente de pares de entrenamiento del sintetizador y, en el único paciente adquirido dos veces con el mismo implante, de par de reproducibilidad. La geometría del implante es paramétrica (Sección~\ref{sec:geometria}), y la auditoría local motiva esa elección. Los tornillos extraídos con un umbral fijo de HU aparecen fragmentados, con diámetros de fuste por debajo de los calibres publicados, y su geometría cambia entre dos adquisiciones del mismo implante.

\section{Evaluación de la representación multiventana}\label{sec:obj1}

La compuerta del Objetivo~1 existe porque un modelo de difusión que opera en el espacio latente de un autoencoder solo puede generar lo que ese autoencoder es capaz de representar. Si el paso de ida y vuelta (HU $\rightarrow$ representación $\rightarrow$ HU) pierde los valores del hueso denso, los HU del hueso quedan sesgados en el volumen sintetizado aunque el generador no tenga error. La compuerta se definió, por eso, como obligatoria, de tipo \emph{Go/No-Go}, y previa al Objetivo~3. Su regla operativa se fijó por escrito antes de correr la prueba que decide la compuerta. Una exploración previa del autoencoder preentrenado sobre los 178 volúmenes locales, incluidos los de los pacientes que después formaron la partición de prueba, ya había quedado por encima del criterio en todos ellos. La regla se fijó, por tanto, conociendo un resultado exploratorio desfavorable.

La cadena evaluada convierte cada corte en una \textbf{codificación multiventana} de tres canales, cada uno con una ventana de HU distinta. Esa codificación pasa por el codificador y el decodificador del autoencoder variacional de Stable Diffusion 1.5, y los HU se recuperan sin usar la imagen original. La lectura toma, en cada vóxel, el canal más estrecho que no esté saturado. El error se mide dentro del hueso, delimitado por el mismo umbral de 150~HU que usa el protocolo de evaluación adoptado \cite{peters2025hybrid,haneda2025aapm}. Para un paciente $i$ con conjunto de vóxeles óseos $\Omega_i$, el error es el error absoluto medio (MAE, \emph{mean absolute error}):
\begin{equation}
\mathrm{MAE}_i = \frac{1}{|\Omega_i|} \sum_{v \in \Omega_i} \left| x_v - \hat{x}_v \right|,
\label{eq:mae}
\end{equation}
donde $x_v$ es el valor original en HU y $\hat{x}_v$ el valor recuperado tras la ida y vuelta.

El umbral de 25~HU no proviene de un criterio publicado, porque no se encontró ningún umbral de aprobación sobre exactitud de los valores de TC. Se ancla en los errores que reportan tres métodos de MAR en la misma escala. Karageorgos et al.~\cite{karageorgos2024ddpm} reportan una raíz del error cuadrático medio (RMSE, \emph{root mean square error}) de 20.2~HU para el algoritmo que ancla la escala del protocolo de Peters et al.~\cite{peters2025hybrid}. Para su modelo de difusión en el dominio del sinograma, evaluado en el dominio de imagen, reportan 12.3~HU, y Yun et al.~\cite{yun2026simulationdriven} reportan 12.74~HU para un modelo de difusión latente. Como para los mismos residuos el RMSE nunca es menor que el MAE, la comparación fija un orden de magnitud y no una equivalencia entre las dos métricas.

El diseño cruza dos variantes del autoencoder con tres codificaciones, lo que da seis combinaciones. Las variantes son el autoencoder preentrenado y una versión con el decodificador adaptado y el codificador congelado, con un decodificador adaptado por cada codificación. Cada decodificador se ajustó con cortes codificados de los pacientes de entrenamiento, separados de los de prueba \GAPDATO{parámetros del ajuste del decodificador (número de cortes, función de pérdida, optimizador y criterio de parada): la pérdida L1 y el optimizador AdamW constan en el código del experimento, y el resto solo como valores por omisión de ese código, no en un registro}. La primera codificación usa las ventanas reportadas por Wang et al.~\cite{wang2025adaptiveweighting}: $[-1000, 2000]$, $[-320, 480]$ y $[-160, 240]$~HU. La segunda eleva el techo de la ventana ancha a 20\,000~HU, y la tercera sustituye la ventana ancha por una compresión arcoseno hiperbólico sobre $[-1000, 20\,000]$~HU.

Una combinación aprueba cuando la media, sobre los 34 pacientes de prueba, del $\mathrm{MAE}_i$ de la Ecuación~\eqref{eq:mae} es menor que 25~HU, y la compuerta aprueba si lo hace alguna de las seis. Esa multiplicidad no se corrige; en su lugar, cada veredicto se acompaña del intervalo de confianza (IC) del 95\,\% por remuestreo de pacientes, y un aprobado cuyo límite superior alcance 25~HU se reporta como marginal. Si aprueban varias combinaciones, la que pasa al Objetivo~3 se elige por un orden fijado a priori y no por su error. Ese orden es arcoseno hiperbólico, luego el techo de 20\,000~HU y por último las ventanas publicadas, cuyo techo de 2000~HU satura el metal; dentro de cada codificación, el decodificador preentrenado precede al adaptado. Si ninguna aprueba, la regla registrada establece que el Objetivo~3 no se ejecuta y que el veredicto negativo se reporta como resultado.

En los volúmenes con metal, el mismo error se reporta además dentro de $B_{\delta}$, de forma descriptiva y fuera de la regla de la compuerta. En esta medición, $B_{\delta}$ son los vóxeles a 12~mm o menos, en distancia euclídea tridimensional, del metal real delimitado por el umbral de 2500~HU, excluido el propio metal. Esa banda contiene las rayas (\emph{streaking}) de baja amplitud que el sintetizador debe generar, y un error restringido al hueso no muestra si el autoencoder las conserva. Rombach et al.~\cite{rombach2022latentdiffusion} diseñan la primera etapa de la difusión latente para eliminar detalle de alta frecuencia, de modo que esa conservación no puede suponerse.

Después del veredicto negativo, la compuerta se extendió a un latente entrenado con TC, el del generador de volúmenes de TC de Guo et al.~\cite{guo2025maisi}, con la misma regla y los mismos pacientes de prueba. Antes de correr ese modelo se comprobó su rango de salida, porque saturar las intensidades en ese rango fija una cota inferior del error que ningún peso puede mejorar. Esa cota se calcula saturando los volúmenes de prueba al rango del modelo, sin modelo alguno, y si ya supera el criterio la evaluación no se corre. El veredicto de la compuerta y el de esta extensión se presentan en el capítulo de resultados.

\section{Muestreador de colocación quirúrgicamente restringido}\label{sec:muestreador}

El muestreador no busca una trayectoria óptima: perturba una trayectoria ya medida sobre la anatomía de cada paciente. La razón es que no existe una pose canónica que reproducir. Zwingmann et al.~\cite{zwingmann2009navigated} reportan distribuciones de malposición distintas para una serie operada con navegación y otra con técnica convencional; esas dos series forman la \textbf{referencia clínica} de este trabajo. Además, la orientación de las superficies óseas que limitan el corredor tuvo, entre especímenes cadavéricos, coeficientes de variación de 7 a 25\,\% en la mayoría de las superficies y de hasta 140\,\% en el ala superior de S1 \cite{ziran2007fluoroscopic}. El muestreador produce entonces una distribución de poses preinscrita (Sección~\ref{sec:poses}). Cada pose recibe un \textbf{grado de brecha cortical}, que mide cuánto sobresale el implante del hueso (Sección~\ref{sec:sap}), y la distribución de grados que resulta se compara con la referencia clínica sin haberse ajustado a ella.

\subsection{Marco de referencia y medición del corredor}\label{sec:corredor}

Las poses se expresan en el marco de referencia de Kaiser et al.~\cite{kaiser2014dysmorphism}, que es computable sobre un volumen de TC. El volumen se reformatea a lo largo del eje sacro, perpendicular al platillo superior de S1. La angulación del corredor se mide en el plano coronal contra la línea que une las crestas ilíacas, y en el plano axial contra la línea que une las espinas ilíacas posteriores. El marco incluye además una regla de longitud útil que exige al menos 5~mm hasta la cortical a cada lado del eje. Este trabajo lee ese \textbf{margen cortical de longitud útil}, $h$, como una distancia radial medida en el plano perpendicular al eje, y no como una holgura añadida al tornillo. Con esa lectura, $h$ coincide con el radio del corredor de 10~mm de Kaiser et al.

La geometría del corredor se mide sobre máscaras anatómicas y no sobre un umbral de HU. A lo largo del eje del corredor más ancho, la mediana de la fracción de vóxeles a 150~HU o menos dentro de las máscaras sacras fue 0.42, con el mismo umbral de hueso del protocolo adoptado (Sección~\ref{sec:obj1}). Una definición de hueso por HU dejaría fuera, en la mediana, esa fracción del corredor.

Las máscaras provienen de TotalSegmentator 2.18.0 \cite{wasserthal2023}, construido sobre nnU-Net \cite{isensee2021}, con la tarea \texttt{total}. Antes de segmentar, la herramienta delimita la región de interés con un modelo de recorte de menor resolución. El \textbf{recorte por defecto} usa un modelo de 6~mm; el \textbf{recorte alternativo}, activado por la opción \texttt{-{}-robust\_crop}, uno de 3~mm. El primero es el principal y el segundo se usa como análisis de sensibilidad. Cada caso pasa el control de nivel descrito en la Sección~\ref{sec:datos}, y se eliminan los componentes conexos menores que el 0.1\,\% de su estructura. La publicación citada evalúa un modelo de 104 estructuras y no reporta exactitud para la etiqueta de S1, por lo que esa exactitud no se asume para la versión y las clases usadas aquí.

La \textbf{envolvente ósea} es la unión de las máscaras de sacro, vértebra S1 y ambos coxales, con un cierre morfológico de 2~mm y sin relleno de cavidades cerradas. La decisión que la definía con ese relleno se corrigió, porque las máscaras de TotalSegmentator son volúmenes sólidos y no dejan hueco que tapar. Aplicar el relleno no cambia el diámetro del corredor en ninguno de los 72 pacientes de la cohorte primaria. El diámetro del corredor se define como el doble de la distancia libre mínima a esa envolvente a lo largo del eje, excluidos 8~mm en cada extremo (Sección~\ref{sec:sap}). Por cada caso, la búsqueda devuelve el eje del corredor más ancho, con su centro, su dirección y su longitud \GAPDATO{descripción completa del algoritmo de búsqueda del eje del corredor; hoy consta en el código y, en parte, en el registro de implicancias del proyecto}. Si en el volumen hay un implante, este se trata como espacio ocupado mediante un umbral de semimáximo local por objeto, y un umbral fijo de 2500~HU da la cota superior del corredor ocupado. Esos volúmenes se reportan como contraste aparte y no se agregan a la cohorte del corredor.

Un corredor es viable cuando su diámetro $D$ admite el calibre del tornillo más una \textbf{holgura radial} a cada lado, $D \geq d + 2\epsilon$, donde $d$ es el calibre nominal del tornillo paramétrico (Tabla~\ref{tab:geometrias}) y $\epsilon$ vale entre 1 y 2~mm. La holgura radial se toma de Kaiser et al.~\cite{kaiser2014dysmorphism}, que justifican su corredor de 10~mm como holgura de 1 a 2~mm alrededor de un tornillo de 6.3 a 8~mm. Leerla como holgura radial por lado es una operacionalización propia, porque la publicación no indica si la holgura se cuenta por lado o en total. Kaiser et al. eligen ese corredor de 10~mm como un valor conservador y lo atribuyen a trabajos previos. McLaren et al.~\cite{mclaren2021corridor} lo toman de Kaiser et al. y lo emplean con un procedimiento reproducible de medición del diámetro, pero no lo validan como umbral de seguridad. El umbral fijo de 10~mm se conserva solo como convención de comparación con esa literatura y no como estándar clínico heredado.

Con el recorte por defecto, el diámetro del corredor tuvo en la cohorte primaria una mediana de 9.5~mm (rango intercuartílico de 7.4 a 11.7~mm). Con ese mismo recorte, el criterio $D \geq d + 2\epsilon$ se cumplió en el 65.3\,\% de los 72 pacientes en su extremo menos exigente ($d = 6.5$~mm, $\epsilon = 1$~mm). En el extremo más exigente ($d = 8.0$~mm, $\epsilon = 2$~mm, Tabla~\ref{tab:geometrias}) se cumplió en el 23.6\,\%. Con el recorte alternativo, esos porcentajes fueron 62.5 y 20.8\,\%. La convención de 10~mm se cumplió en 29 de 72 pacientes con el recorte por defecto y en 27 con el alternativo. La diferencia entre recortes es incertidumbre de la medición, no variabilidad anatómica. La limpieza de componentes no cambió el diámetro ni su viabilidad en ninguna fracción probada hasta el 5\,\%; la fracción de 0.1\,\% que se usa, como la regla de ocupación, se fijó después de inspeccionar estos resultados.

Las cifras principales del Objetivo~2 usan el recorte por defecto, y la decisión que lo fijó quedó condicionada a revisar los 16 casos marcados por diferencia de diámetro entre recortes o por desplazamiento de la caja de una estructura. Esa revisión se hizo sobre láminas, por la autora con apoyo de un médico egresado y sin especialidad, y cubrió los 16 casos. De ellos, 14 salieron conformes, 2 fallaron con los dos recortes y ninguno falló solo con el de 6~mm, que es el único veredicto que habría reabierto la decisión. El recorte por defecto se mantiene, ahora con la condición evaluada y no solo declarada; los 2 casos que fallan con ambos recortes son errores de segmentación y no de recorte \GAPDEC{tratamiento de los dos casos con error de segmentación anotado en la revisión de láminas}.

\subsection{Aplicabilidad del marco en volúmenes con metal}\label{sec:marco-metal}

El marco de Kaiser et al. se apoya en referencias anatómicas de la unión lumbosacra, que es donde se concentran las rayas del artefacto metálico. Ninguna de las fuentes revisadas dice si esas referencias siguen siendo localizables con metal, de modo que se cuantificó en la cohorte local. Las referencias se localizaron con una heurística automática \GAPDATO{descripción de la heurística de localización de las referencias anatómicas, que hoy consta en el código y, en parte, en las notas del experimento}, y una región se contó como contaminada cuando su fracción de rayas oscuras o de metal superaba el máximo observado en 69 pacientes de \texttt{dataset6} que la revisión tridimensional clasificó sin objeto. Ese conjunto de calibración no coincide con el grupo~3 de la Sección~\ref{sec:datos}, que tras el reparto por grupos reúne 66 pacientes. Un revisor clínico, cirujano otorrinolaringólogo y ciego a la auditoría automática, revisó el nivel S1 en cada paciente.

De 65 pacientes con material ortopédico, 57 tenían las cinco referencias localizadas y dentro del campo de visión. De los 8 restantes, 7 tenían las crestas ilíacas truncadas por el campo de visión. El marco fue computable con el nivel S1 correcto según el revisor clínico en 48 pacientes, y quedó libre de contaminación en las cinco regiones en 29. En los 17 pacientes en que el censo morfológico del metal identificó un candidato a tornillo transilíaco-transsacro (Sección~\ref{sec:geometria}), el nivel S1 fue correcto en todos y su región estaba contaminada en 3. La localización heurística de crestas y espinas ilíacas no se ha revisado clínicamente.

La autora volvió a juzgar el nivel S1 sobre las mismas 61 láminas sagitales. El acuerdo con el revisor clínico fue de 61 de 61 casos, con dos límites: el segundo lector es la autora y no un segundo clínico, y solo los juicios de nivel son independientes. Las notas de texto libre se copiaron una vez comprobada la coincidencia, de modo que no aportan evidencia independiente. Las revisiones de láminas registran, además, qué estructura se acepta como S1 cuando la máscara se extiende por detrás del cuerpo vertebral. En 12 de los 16 casos revisados por el recorte (Sección~\ref{sec:corredor}) es la cresta sacra media, que la etiqueta \texttt{vertebrae\_S1} de TotalSegmentator contiene por diseño, porque etiqueta la vértebra completa y no solo su cuerpo.

\subsection{Geometría del implante}\label{sec:geometria}

Los implantes se modelan como tornillos rígidos paramétricos en lugar de extraerse de los volúmenes con metal. La trayectoria que este trabajo mide y muestrea va de la cortical externa de un ilion, cruzando las dos articulaciones sacroilíacas, a la cortical externa del ilion contralateral. El implante que representa es, por tanto, un tornillo transilíaco-transsacro, y le aplican las tolerancias angulares de McLaren et al.~\cite{mclaren2021corridor}, medidas sobre ese mismo corredor. Se mantienen separadas tres geometrías, una para decidir la viabilidad, otra para sintetizar y otra para medir la brecha, que la Tabla~\ref{tab:geometrias} resume. La cifra nominal de estos tornillos es el diámetro de la rosca, no el del cuerpo. Un cilindro uniforme de 6.5 a 8.0~mm tendría, a lo largo del corredor, más sección de metal que el cuerpo de 4.91~mm de la Tabla~\ref{tab:geometrias}. El calibre nominal solo se usa para la viabilidad, y la máscara que se sintetiza es el cuerpo del tornillo.

\begin{table}[htbp]
\centering
\small
\caption{Las tres geometrías del implante y su propósito.}
\label{tab:geometrias}
\begin{tabular}{p{0.2\linewidth} p{0.2\linewidth} p{0.5\linewidth}}
\hline
\textbf{Geometría} & \textbf{Propósito} & \textbf{Valor y respaldo} \\
\hline
Calibre nominal & Viabilidad del corredor & 6.5 a 8.0~mm \cite{gardner2010safezones}; 6.3 a 8~mm \cite{kaiser2014dysmorphism} \\
Máscara del implante $M$ & Lo que se sintetiza (Objetivo~3) & Cilindro uniforme de 4.91~mm, mediana del diámetro del cuerpo medida en la cohorte; fuste de 4.8~mm en la guía del fabricante \cite{synthes2003guide}; 4.9~mm de fuste y 4.7~mm de núcleo bajo la rosca \cite{gardner2015screw} \\
Calibre de medición & Lo que mide la brecha en SAP & 7.0~mm, el de los tornillos canulados de la referencia clínica \cite{zwingmann2009navigated}; 4.91 y 7.3~mm como sensibilidad \\
\hline
\end{tabular}
\end{table}

El cuerpo de 4.91~mm es la mediana medida en la cohorte local sobre 79 componentes metálicos alargados y finos de 43 volúmenes. Los seleccionó el \textbf{censo morfológico del metal}, un filtro geométrico y no clínico: vóxeles sobre 2500~HU, longitud de al menos 30~mm y anchos de 12~mm o menos. Ese censo no afirma que los componentes sean tornillos transilíaco-transsacros, sino que son candidatos por su forma. La cifra está próxima al fuste de 4.8~mm que imprimen la guía técnica del fabricante \cite{synthes2003guide} y el modelo de elementos finitos de Zhu y Liao~\cite{zhu2022optimalposition}. La guía da ese mismo fuste para los tornillos de 6.5 y 7.3~mm, y el modelo usa 7.3~mm solo en una rosca de 16~mm.

La cabeza, cuando se modela, mide 8.0~mm de diámetro y 4.5~mm de altura. El diámetro lo reportan Sayres et al.~\cite{sayres2014comparison} y lo imprime también el catálogo de otro fabricante \cite{doublemedical2021trauma}; la altura, entre las fuentes revisadas, solo la dan Sayres et al. Que la cabeza quede sin avellanar es un supuesto propio; la cabeza y la rosca entran como variantes de sensibilidad.

El calibre de medición lo fija la referencia clínica y no se elige. Zwingmann et al.~\cite{zwingmann2009navigated} asignaron sus cuatro grados sobre tornillos canulados de 7.0~mm, y la profundidad de perforación crece con el radio. Medir las colocaciones simuladas con otro calibre introduciría en la distancia de comparación un sesgo sistemático que no se podría separar del error de colocación. La longitud del tornillo tampoco procede de un rango publicado: se acota en cada volumen por el corredor medido, con la regla de longitud útil de Kaiser et al.

Tres parámetros quedan sin fijar por la documentación del propio fabricante y se declaran en lugar de suponerse. El espesor de la arandela se toma como 1.5~mm del catálogo de otro fabricante \cite{doublemedical2021trauma}. Las dos fuentes de fabricante \cite{synthes2003guide,doublemedical2021trauma} son documentación técnica no revisada por pares. El diámetro de la canulación queda como parámetro libre y se evalúa el tornillo macizo frente al hueco, porque el valor que circula en listados de distribuidores solo aparece en la guía del fabricante como canulación de brocas y destornilladores. La aleación tampoco se codifica: el mismo sistema se fabrica en acero inoxidable 316L y en Ti-6Al-7Nb \cite{synthes2003guide}. Una máscara binaria no distingue entre ambos, de modo que el sintetizador se entrena con la apariencia de una mezcla de materiales.

\subsection{Distribución de poses preinscrita}\label{sec:poses}

La distribución de poses y la métrica se congelaron por escrito el 22 de septiembre de 2026, antes de calcular cualquier distancia contra la referencia clínica. Esa \textbf{preinscripción} (\emph{pre-registration}) impide que la comparación sea un ajuste: si la distribución se hubiera modificado después de ver la distancia, se compararía el modelo contra los mismos datos usados para construirlo. Cualquier cambio posterior se registra como desviación, con fecha y motivo.

Para cada caso, la pose base es el eje del corredor medido en ese volumen, anclado en el punto medio de su tramo de corredor, y de ese mismo tramo sale la longitud del implante. El anclaje corrige un desajuste de la búsqueda: el punto desde el que se lanza la búsqueda no siempre coincide con el centro del tramo óseo encontrado. Sin él, un implante simétrico alrededor de ese punto quedaría descentrado, y la identidad ``implante sobre el eje del corredor, brecha nula'' dejaría de cumplirse.

Sea $\mathbf{c}$ el centro anclado y $\mathbf{u}$ la dirección unitaria del eje. Cada pose aplica dos perturbaciones en el plano perpendicular a $\mathbf{u}$, ambas con dirección uniforme en ese plano y magnitud seminormal truncada en tres desviaciones estándar:
\begin{equation}
\mathbf{c}' = \mathbf{c} + \rho\,\mathbf{e}_{\phi}, \qquad \rho = |N(0, \sigma_t)|, \qquad \mathbf{u}' = R_{\psi}(\alpha)\,\mathbf{u}, \qquad \alpha = |N(0, \sigma_a)|,
\label{eq:pose}
\end{equation}
donde $\mathbf{e}_{\phi}$ es un vector unitario perpendicular a $\mathbf{u}$ con ángulo $\phi$ uniforme, y $R_{\psi}(\alpha)$ es una rotación de ángulo $\alpha$ alrededor de un eje perpendicular a $\mathbf{u}$ de orientación $\psi$ uniforme. La Tabla~\ref{tab:preinscripcion} recoge los parámetros.

\begin{table}[htbp]
\centering
\small
\caption{Parámetros preinscritos de la distribución de poses.}
\label{tab:preinscripcion}
\begin{tabular}{p{0.3\linewidth} p{0.22\linewidth} p{0.38\linewidth}}
\hline
\textbf{Parámetro} & \textbf{Valor} & \textbf{Procedencia} \\
\hline
Margen cortical $h$ & 5.0~mm & Regla de longitud útil de Kaiser et al.~\cite{kaiser2014dysmorphism} \\
$h = 2\sigma$ & convención & Propia y declarada; no es un dato medido \\
$\sigma_t$ & $h/2 = 2.5$~mm & Derivado \\
$\sigma_a$ & $\frac{1}{2}\arctan\!\left(\frac{h}{L/2}\right)$, por caso & Derivado; $2.07^{\circ}$ con la mediana de longitud del corredor, $L = 138$~mm \\
Truncamiento & $3\sigma$ & Declarado; acota la magnitud de cada perturbación \\
Poses por caso & 50 & Declarado \\
Semilla & 20260922, derivada por caso & Fija; las poses no se vuelven a sortear \\
\hline
\end{tabular}
\end{table}

Toda la escala deriva de una sola constante independiente, el margen cortical de 5~mm de la regla de longitud útil (Sección~\ref{sec:corredor}). La convención propia fija ese margen en dos desviaciones estándar: en distancia para $\sigma_t$ y en el ángulo que produce para $\sigma_a$ \GAPDEC{justificación de la convención $h = 2\sigma$, que fija la escala de todas las perturbaciones y no consta en el repositorio}. El margen $h$ es una distancia radial al eje, es decir, se mide en el mismo plano perpendicular en que actúan las perturbaciones. La dispersión angular $\sigma_a$ es ese mismo margen expresado como el ángulo que lo produce en la punta del corredor de cada caso. Ningún parámetro se toma de Zwingmann et al.~\cite{zwingmann2009navigated}, incluida la tolerancia angular de $4^{\circ}$ que su introducción cita de una tercera fuente: aparece en el mismo artículo contra el que se puntúa el resultado. La perturbación es isótropa porque privilegiar una dirección anatómica de error exigiría un dato clínico de direccionalidad que no se tiene. Fijarla a juicio añadiría un parámetro sin dato que lo sostenga.

La semilla se deriva del identificador de cada caso. La cohorte de sensibilidad es un subconjunto de la primaria, y con un generador que avanzara de caso en caso el mismo paciente recibiría poses distintas en las dos corridas. Con la semilla por caso, cada paciente recibe las mismas poses en cualquier cohorte que lo contenga, y las dos cohortes difieren solo en su composición.

Antes de la corrida se fijaron tres comprobaciones que la implementación debía cumplir. Son controles de consistencia de la implementación y no condiciones de fallo del muestreador como modelo; las dos primeras forman parte de los seis controles de SAP (Sección~\ref{sec:sap}), y la tercera se comprueba en cada corrida. La pose sin perturbar debe dar brecha nula con el diámetro de su propio corredor, porque es el eje del corredor. La brecha no debe decrecer al ensanchar el cilindro. Ninguna pose puede quedar sin calificar, porque descartar las que no atraviesan hueso eliminaría del denominador las peores colocaciones y bajaría la distancia reportada sin que la colocación mejorara.

Bajo el corredor de S1, la búsqueda devuelve un segundo corredor óseo cuyo nivel vertebral no es constante. En una submuestra de 18 volúmenes estratificada por profundidad, una segunda lectura de la autora, ciega a la profundidad medida, lo ubicó en S2 en 13 volúmenes, en S3 en 4 y en S4 en 1. La profundidad no determina el nivel: un corredor a 39~mm bajo S1 se juzgó S2, mientras otros a 36 y 33~mm se juzgaron S3. De ahí que se nombre ``segundo corredor bajo S1'' y no S2. Sus resultados son descriptivos, porque ninguna de las fuentes revisadas reporta una distribución ordinal clínica por debajo de S1, y la preinscripción cubre solo S1 \GAPDATO{preinscripción y corrida del muestreo de poses en el segundo corredor bajo S1; su eje ya está medido en 69 de los 72 pacientes de la cohorte primaria}.

La densidad ósea entra como componente reportado de SAP y no como condicionante del muestreo. Si condicionara el muestreo, se volvería un parámetro ajustable más y reabriría la circularidad que la preinscripción cierra. La motivación es la heterogeneidad de la masa ósea pélvica que describen Arand et al.~\cite{arand2019pelvicring} en su modelo estadístico poblacional, por ejemplo los valores de gris bajos en el ala sacra \GAPDEC{fuente y definición operativa de la fracción por zona de densidad: la decisión que la fija como componente de SAP la calcula sobre una sonda de densidad que una decisión anterior retiró como evidencia de densidad, y la definición implementada solo consta en el código}.

Los fenotipos de morfología sacra no se usan como variable de estratificación, porque su efecto sobre el tamaño de la zona segura no es consistente entre estudios. Gardner et al.~\cite{gardner2010safezones} encuentran una zona segura de S1 menor en sacros dismórficos, pero recogen un estudio previo que no halló diferencia entre sacros normales y dismórficos. El muestreador mide, en cambio, la geometría del corredor directamente en cada volumen.

\section{Sintetizador condicionado en el dominio de imagen}\label{sec:sintetizador}

El sintetizador se formula como \emph{inpainting} de la región de generación $G = M \cup B_{\delta}$ mediante un modelo de difusión condicional 2.5D en el dominio de imagen, es decir, sobre valores de TC por vóxel, sin espacio latente ni datos de proyección. El modelo recibe el parche con $G$ borrada, la máscara del implante $M$ y la máscara de la banda, y genera valores de TC solo dentro de $G$. Todo lo que queda fuera de $G$ se copia del volumen de origen, de modo que la preservación fuera de la banda se cumple por construcción y no por medición. Las entradas y salidas usan la codificación multiventana como representación de canales, sin la etapa de compresión del autoencoder. La codificación elegida entre las tres de la compuerta, la arquitectura de la red y los parámetros de entrenamiento y muestreo constan en un diseño que todavía no se congeló \GAPDEC{congelar la preinscripción del diseño del sintetizador: codificación multiventana (el borrador supone la de arcoseno hiperbólico), arquitectura, los demás parámetros de entrenamiento, muestreo, número de semillas por caso y margen de equivalencia}.

La codificación multiventana acota por abajo los valores de TC representables, porque su ventana ancha empieza en $-1000$~HU (Sección~\ref{sec:obj1}). Qué se hace con ese suelo lo decide una regla preinscrita. La regla mide, sin modelo y en la banda $B_{\delta}$ de los pacientes con metal de entrenamiento y validación, qué parte de la diferencia entre las colas del 5\,\% de los HU pierde la ida y vuelta por la representación. Esa pérdida se evalúa en el percentil~5 de cada paciente y no en la mediana, una elección hecha después de ver su distribución. Si es menor que el margen de equivalencia $\Delta$ (Sección~\ref{sec:apariencia}), el suelo se declara como limitación de alcance; si es mayor, se cambia la representación \GAPDATO{desenlace de la regla preinscrita sobre el suelo de $-1000$~HU, que depende del margen de equivalencia $\Delta$, aún no medido}.

Esta formulación sustituye a la inicial, que era difusión latente sobre la base preentrenada Stable Diffusion 1.5 guiada con ControlNet. ControlNet presupone una base preentrenada congelada y su espacio latente \cite{zhang2023controlnet}, y el Objetivo~1 midió que ese latente no conserva los valores de TC del hueso con la tolerancia exigida. Como no se adopta ninguna base de TC preentrenada en el dominio de imagen, la guía por ControlNet pierde su objeto.

Los HU de la salida del sintetizador se leen de su codificación multiventana con una regla propia del Objetivo~3, distinta de la del Objetivo~1. La lectura del Objetivo~1 toma el canal más estrecho que no esté saturado (Sección~\ref{sec:obj1}), y con ella un canal estrecho casi saturado sustituía el valor del canal de ventana ancha por el techo de su propia ventana. El sintetizador parte, en cambio, del canal de ventana ancha y lo corrige con cada canal estrecho mediante una mezcla ponderada, en la que un canal pesa menos cuanto más se acerca a la saturación. El peso depende de la distancia del valor normalizado del canal al borde de su ventana: es nulo por debajo de 0.01, crece de forma lineal y es completo desde 0.05. Si los dos límites coinciden en 0.01, que es el umbral de saturación de la lectura del Objetivo~1, la rampa se reduce a un escalón y la mezcla da esa misma lectura, que queda así como caso particular. El límite de 0.05 se fijó sobre el conjunto de validación, con un criterio registrado antes de conocer el resultado \GAPDATO{resultado citable de la calibración de ese límite sobre el conjunto de validación}. La compuerta del Objetivo~1 conserva su lectura, y sus cifras no cambian.

La banda $B_{\delta}$ extiende la síntesis unos 12~mm más allá del borde del implante, para que las rayas y el endurecimiento del haz puedan manifestarse fuera de la máscara. Su ancho es un parámetro de diseño: acota dónde se genera artefacto y trunca a propósito las rayas lejanas, cuya extensión espacial no está medida en la literatura revisada. Ninguna de las fuentes revisadas da un valor para el ancho. Karageorgos et al.~\cite{karageorgos2024ddpm} sí dan el sentido del error en su propio dominio: estudian la sensibilidad de su método a la máscara metálica, con la que obtienen un RMSE de 7.57~HU cuando la máscara es la verdadera. Dilatarla eleva el error a 11.45 y 11.63~HU para factores de 1.4 y 1.8; contraerla lo eleva a 54.82 y 57.69~HU para factores de 0.7 y 0.5. Ese experimento varía, sin embargo, una máscara de reducción de artefactos en el dominio del sinograma y no una banda de generación en el dominio de imagen. Este trabajo lo lee como apoyo a generar más allá de la máscara del implante, pero no como valor para el ancho de 12~mm.

Los pares de entrenamiento se forman con los volúmenes con metal de los pacientes de entrenamiento, donde el implante real aporta a la vez la máscara y el objetivo de apariencia. La máscara de entrenamiento se obtiene con el umbral de 2500~HU, y el semimáximo local por objeto entra como sensibilidad. La unidad de entrenamiento es cada componente conexo de metal y no el corte axial, para que una región de entrenamiento tenga la misma estructura que una región de síntesis, que contiene un solo implante. En la síntesis, la máscara es el tornillo paramétrico de la Sección~\ref{sec:geometria} colocado en una pose muestreada. El entrenamiento se fija en 30\,000 pasos, una elección de hiperparámetro hecha sobre el conjunto de validación y no fijada de antemano. Dos corridas sitúan el mínimo de la curva de validación en una meseta, y ese valor ocupa su centro. Ningún paciente de prueba intervino en la elección, de modo que la regla de aislamiento estricto (Sección~\ref{sec:datos}) se mantiene.

La inclusión en el conjunto de entrenamiento se rige por un criterio fijado antes de entrenar, con tres reglas. El componente debe estar dentro del cuerpo del paciente, y el caso debe tener material ortopédico declarado en la auditoría. El parche, por último, debe contener metal del componente objetivo. A esos parches se añaden todos los parches disponibles que solo contienen banda, que están en razón de 0.62 respecto de los parches con metal. El tamaño y la forma del componente se reportan y no se filtran, porque se asume que la relación entre máscara y artefacto no depende del tipo de implante. Ese supuesto no se verificó. Lo que sí se midió es el costo de filtrar por forma: en una muestra revisada de 40 componentes, solo 5 eran tornillos, de modo que un filtro de forma habría descartado la mayor parte de los ejemplos.

Entre entrenamiento y síntesis hay tres desplazamientos de dominio. Primero, las máscaras de entrenamiento se obtienen umbralizando metal real, cuyo fuste la auditoría midió por debajo de los calibres publicados, mientras que la máscara de síntesis es un cilindro paramétrico liso. Segundo, fuera de $G$ el contexto trae rayas lejanas en los volúmenes de entrenamiento con metal, pero no en los volúmenes sin metal que reciben la síntesis. Tercero, los parches que solo contienen banda son menos frecuentes en el conjunto de entrenamiento de lo que la geometría del corredor predice para la síntesis: 0.62 frente a 1.40 parches de solo banda por parche con metal. El tercero queda así cuantificado. Para el primero falta la métrica \GAPDEC{métrica y análisis con que se cuantifica la diferencia entre la máscara umbralizada de entrenamiento y el cilindro liso de síntesis}. Para el segundo, el salto de los valores de TC al cruzar el borde de la banda ya se midió una vez, sobre un caso, un corte y un modelo sin converger \GAPDEC{preinscribir el bloque que mide la continuidad de los valores de TC en el borde de la banda y su análisis}.

El \textbf{punto de control} (\emph{checkpoint}) del sintetizador que pasa a la evaluación no se elige por la pérdida de validación, porque esa pérdida no ordena los puntos de control por la fidelidad de los HU que generan. Esa misma razón debilita el argumento de los 30\,000 pasos, tomado del mínimo de esa curva: la cifra se mantiene y se revisará cuando el criterio de selección esté medido. Tampoco se elige por la reconstrucción de implantes reales de validación, porque esa tarea premia memorizar implantes parecidos a los de entrenamiento y no es la de este trabajo, que coloca un tornillo donde no había metal. El criterio se mide, por eso, sobre la tarea de síntesis. Se coloca el tornillo paramétrico en una pelvis de validación sin osteosíntesis, se genera con cada punto de control candidato y se comparan dos estadísticos de la apariencia sintética con los mismos estadísticos medidos en los implantes reales de validación. Ninguno de los dos exige una imagen sin metal del paciente, que en un implante real no existe y que la \emph{streak amplitude} sí necesita (Sección~\ref{sec:apariencia}). Los HU sintéticos se leen con la mezcla de canales descrita arriba.

El primer estadístico es el perfil radial de HU alrededor del metal. Se mide en el plano de cada corte, por cáscaras de 0.5~mm de 0 a 12~mm en anillo completo. En cada cáscara se toman la mediana y el percentil~95 de HU, expresados como elevación sobre la mediana del anillo de 12 a 15~mm del mismo paciente. El percentil~95 acompaña a la mediana porque las rayas ocupan las colas de la distribución de HU, que la mediana no registra. Los valores se agregan por corte, después por paciente y después entre pacientes. El segundo estadístico es el histograma de HU dentro de $M$, resumido por paciente en los percentiles~50, 75 y~95 de los vóxeles sobre 2500~HU. La fracción sintética de $M$ que supera 2500~HU se reporta solo como control y no se compara con la real, porque en los implantes reales la máscara se define con ese mismo umbral y esa fracción no informa nada.

Como referencia se usan los implantes reales de los 3 pacientes de validación (Sección~\ref{sec:apariencia}). Gana el punto de control con más cáscaras dentro de la envolvente real, del mínimo al máximo de esos 3 pacientes, en las dos curvas del perfil. Un empate lo decide la menor distancia a la mediana real de los percentiles~50 y~95 dentro de $M$, y un empate que persista lo decide la autora. Con 3 pacientes de referencia, el cotejo puede descartar un punto de control cuyo artefacto cae fuera de lo observado, pero no prueba que el elegido lo reproduzca. Exige además que esos 3 pacientes lleven tornillos, y su tipo de implante no está verificado; si no lo son, el cotejo se revisa \GAPDEC{verificación del tipo de implante de los 3 pacientes de referencia antes del cotejo; el procedimiento para verificarlo es una propuesta preliminar}. Mientras el cotejo no esté medido, ningún punto de control está elegido.

\section{Protocolo de evaluación}\label{sec:evaluacion}

La evaluación separa la colocación de la apariencia porque cada una tiene su propia referencia. La colocación se juzga contra la referencia clínica; la apariencia, contra la \textbf{inserción por copia y pegado} y contra el protocolo físico de Peters et al. La Tabla~\ref{tab:diseno} resume el diseño experimental de los tres objetivos que se evalúan con datos. El Objetivo~4 no tiene fila propia: su producto es la definición de las métricas que usan las filas~2 y~3. La implementación de SAP se verificó con los controles de la Sección~\ref{sec:sap} \GAPDEC{qué evidencia verifica el Objetivo~4 más allá de esos controles, o si se declara que no se evalúa experimentalmente}.

\subsection{Admisibilidad quirúrgica de la colocación}\label{sec:sap}

SAP tiene tres componentes: el grado de brecha cortical de cada pose, la fracción por zona de densidad y la viabilidad del corredor en el marco de Kaiser et al. Solo la distribución de grados se compara con la referencia clínica. La fracción por zona de densidad y la viabilidad se reportan de forma descriptiva, la segunda por caso y con el criterio $D \geq d + 2\epsilon$ de la Sección~\ref{sec:corredor} \GAPDEC{calibre con que SAP evalúa la viabilidad del corredor: la decisión que define SAP fija el calibre nominal de 6.5 a 8.0~mm, y la corrida del Objetivo~2 usó los calibres de 4.91, 7.0 y 7.3~mm con una holgura de 1~mm}. El grado de brecha cortical usa la escala ordinal de cuatro grados que Smith et al.~\cite{smith2006iliosacral} aplican a la fijación iliosacra. Smith et al. suman además una escala angular a la de perforación \GAPDEC{si SAP incorpora la dimensión angular de la escala de Smith et al. o declara que usa solo la de perforación, y con qué justificación}. La brecha de una pose es la protrusión radial máxima del implante fuera de la envolvente ósea segmentada:
\begin{equation}
b = \max_{t \in T} \; \max\bigl(0,\; r + s(\mathbf{c}' + t\,\mathbf{u}')\bigr),
\label{eq:brecha}
\end{equation}
donde $r$ es el radio del calibre de medición, $s$ es la distancia con signo a la envolvente ósea (negativa dentro del hueso) y $T$ es el tramo evaluado a lo largo del eje. La envolvente es una máscara de hueso y no una segmentación de la cortical, así que su borde aproxima la superficie cortical externa. Los límites de grado se fijan por convención: grado~0 si $b = 0$; grado~1 si $b \in (0, 2)$~mm; grado~2 si $b \in [2, 4]$~mm; y grado~3 si $b > 4$~mm.

El tramo $T$ es el propio implante, de longitud fija igual a la del corredor del caso, centrado en la pose y sin los 8~mm de cada extremo. Esa exclusión de extremos existe porque un tornillo transilíaco-transsacro entra y sale por la cortical del ilion por diseño; sin ella, toda trayectoria válida puntuaría grado~3 por sus propios puntos de entrada y salida. Los 8~mm son los mismos que se excluyen al medir el diámetro del corredor (Sección~\ref{sec:corredor}) \GAPDEC{justificación del valor de 8~mm de la exclusión de extremos, que no consta en el repositorio}. La longitud no depende de dónde esté el hueso, porque si dependiera la métrica premiaría a las poses que se salen del corredor, cuyo tramo se acortaría con ellas. Las poses que no atraviesan hueso reciben grado~3 en lugar de descartarse, para que el denominador coincida con el de la referencia clínica, que calificó todos sus tornillos.

Seis controles verificaron la implementación antes de la corrida, entre ellos la brecha nula sobre el eje y la monotonía en el diámetro. Uno se hizo sobre un fantoma de geometría conocida y cinco sobre el eje del corredor de dos volúmenes reales. La resolución de la métrica en el fantoma es de 0.083~mm.

La distribución de grados de las poses se compara con las dos distribuciones ordinales de Zwingmann et al.~\cite{zwingmann2009navigated}, condicionadas por técnica quirúrgica y en S1, nivel que la serie convencional declara y que en la navegada se asume (Sección~\ref{sec:amenazas}). En la serie navegada, los grados~0 a~3 reúnen el 69, 15, 8 y 8\,\% de los tornillos; en la convencional, el 40, 37, 11.5 y 11.5\,\%. Con los grados tratados como equiespaciados por convención propia, la distancia de Wasserstein-1 entre dos distribuciones $P$ y $Q$ sobre los grados $0$ a $3$ es
\begin{equation}
W_1(P, Q) = \sum_{k=0}^{2} \left| F_P(k) - F_Q(k) \right|,
\label{eq:w1}
\end{equation}
donde $F_P$ y $F_Q$ son las funciones de distribución acumulada. Su unidad es el grado, no el milímetro.

Una distancia grande no se corrige moviendo parámetros del muestreador, que no se ajustó a esas distribuciones. El Objetivo~2 no tiene regla de decisión: ningún valor de la distancia se fijó de antemano como fallo del muestreador \GAPDEC{si se fija una escala de lectura de la distancia de Wasserstein-1, o qué resultado contaría como fallo del muestreador}. La comparación enfrenta una población simulada sobre pelvis receptoras con dos series clínicas pequeñas, de 26 tornillos en 24 pacientes y de 35 tornillos en 32 pacientes \cite{zwingmann2009navigated}; no es una prueba de equivalencia. Los grados se cuentan por tornillo. Los recuentos implican que algunos pacientes recibieron más de uno, aunque el artículo declara un solo tornillo por paciente.

Además del análisis preinscrito, se añade un análisis post hoc estratificado por si el diámetro del corredor admite el calibre de medición de 7.0~mm (Tabla~\ref{tab:geometrias}). En un corredor más estrecho que ese calibre, el propio eje del corredor ya perfora, y el grado~0 es inalcanzable por la anatomía receptora tal como se presenta, no por la colocación. Eso ocurre en 15 de los 72 pacientes de la cohorte primaria. Qué estrecha esos corredores no está establecido. El cribado de fractura halló fractura sacra desplazada en 7 de los 15, un indicio sugerente y no concluyente (Sección~\ref{sec:datos}). Una fracción no cuantificable del estrechamiento puede ser, por tanto, patología y no variabilidad anatómica normal. La estratificación se reporta junto a la cifra preinscrita y no la sustituye. Restringir la cohorte a los corredores que admiten el calibre de 7.0~mm se consideró y se rechazó por dos razones. La primera es que seleccionar pacientes por una variable correlacionada con el resultado subiría la fracción de grado~0 por selección y no por colocación, el mismo sesgo que introduciría ajustar el muestreador a la referencia clínica. La segunda, que Zwingmann et al. no reportan haber excluido pacientes de corredor estrecho, de modo que una cohorte restringida dejaría de ser comparable con la suya.

\subsection{Coherencia de apariencia}\label{sec:apariencia}

La apariencia se evalúa con las métricas del protocolo de Peters et al.~\cite{peters2025hybrid}, con sus nombres publicados: \emph{bone integrity}, \emph{metal integrity} y \emph{streak amplitude}. La \emph{bone integrity} compara el hueso, definido como los vóxeles sobre 150~HU fuera de la geometría metálica, mediante el cambio de volumen y el coeficiente de Sørensen-Dice. La \emph{metal integrity} compara de forma análoga la máscara metálica obtenida de la imagen, con un umbral adaptativo por región. La \emph{streak amplitude} se mide en regiones perpendiculares a las rayas como la diferencia entre el promedio del 5\,\% de desviaciones más altas y el del 5\,\% más bajas respecto de la imagen sin metal del mismo caso. Peters et al. eligen esas regiones de medición sobre las rayas fuertes de las imágenes sin corregir, y declaran como limitación que la elección se hiciera sobre imágenes sin reducción de artefactos, que el propio algoritmo puede alterar \cite{peters2025hybrid}.

Estas métricas se diseñaron para evaluar MAR, no síntesis, y su uso aquí exige invertirlas. La escala de 0 a 4 del protocolo no se traslada, porque se ancla en un algoritmo de MAR (Sección~\ref{sec:obj1}) que no tiene análogo en síntesis. La inversión está definida para la \emph{streak amplitude}, que es el criterio primario, y no para las dos métricas descriptivas \GAPDEC{definición operativa de la inversión de la \emph{bone integrity} y de la \emph{metal integrity} para evaluar síntesis}.

La definición de la \emph{streak amplitude} fija en qué pacientes puede medirse el criterio primario. Su verdad de terreno es la TC original sin metal del mismo paciente, y el campo de desviación es la diferencia en HU entre la imagen sintética y esa original. La partición de prueba reúne 34 pacientes (Sección~\ref{sec:datos}): 20 con implante real y 14 sin metal. La amplitud solo puede medirse en esos 14, porque de un paciente con implante real no hay imagen sin metal. En los 20 restantes lo que se mide es la discrepancia frente a su TC real, con el error absoluto medio en HU dentro de $G$, y se reporta como medida descriptiva.

Las regiones de medición se derivan de la pose y no se trazan a mano. En planos perpendiculares al eje del implante se toman anillos completos, de 360 grados, que van desde una guarda de un vóxel en plano alrededor de $M$ hasta los 12~mm de la banda de generación extendida. Los planos son todos los cortes con máscara del implante menos los 8~mm de cada extremo, la misma exclusión de extremos con que se mide el corredor y el tramo $T$ de SAP (Sección~\ref{sec:sap}). Las regiones quedan así dentro de la banda, que es donde el sintetizador escribe: fuera de ella su amplitud sería nula por construcción. Un estadístico de colas no necesita orientar la región hacia la raya, porque el 5\,\% superior y el 5\,\% inferior encuentran las rayas claras y las oscuras en cualquier punto del anillo. Esta geometría es una desviación declarada del protocolo publicado y no hereda su validación; a cambio, no depende de un lector y es idéntica para los brazos que se comparan.

La guarda se fija en el mínimo por la dirección del sesgo que introduce. Alrededor del metal, el volumen parcial y el artefacto de campo cercano son la misma señal, y separarlos exigiría una imagen sin metal del mismo paciente, que es justo lo que falta donde hay un implante real. Una guarda mayor descartaría señal de artefacto y sesgaría la amplitud a la baja, de modo que se fija en un vóxel en plano, el ancho mínimo que puede ocupar el volumen parcial. Junto a la amplitud se reporta la fracción de vóxeles de la región de medición que quedan exactamente en el suelo de $-1000$~HU de la representación (Sección~\ref{sec:sintetizador}). Si esa fracción es apreciable, la cola inferior queda censurada y la amplitud medida es una cota inferior, con el mismo sentido de sesgo que la guarda.

El realismo se juzga contra las observaciones reales de artefacto (Sección~\ref{sec:datos}), donde la verdad de terreno sin metal no existe. La amplitud de comparación se calcula allí con la forma del estadístico sobre los HU de la banda, y no sobre un campo de desviación, así que no es la misma magnitud que el criterio primario \GAPDEC{qué distancia entre distribuciones de amplitud se usa frente a las observaciones reales de artefacto, y con qué análisis}.

La \emph{streak amplitude} es el único criterio primario, declarado antes de examinar la partición de prueba; las demás métricas son descriptivas. La elección responde a que las rayas son la magnitud que el sintetizador tiene por objeto generar. Fijar un solo criterio primario evita el patrón de comparación múltiple en que basta con que apruebe una de varias métricas. Toda la evaluación del Objetivo~3 es métrica, a diferencia de la del Objetivo~2, que incorpora lecturas humanas del nivel vertebral y del corredor (Secciones~\ref{sec:corredor} y~\ref{sec:marco-metal}) \GAPDEC{si el Objetivo~3 se evalúa solo con métricas, y con qué justificación, o si se añade una lectura humana de la apariencia}.

Sobre ese criterio primario se ordenan tres contrastes, cada uno con lo que puede probar. La comparación de equivalencia frente al protocolo físico es el contraste primario, porque es el único que enfrenta la amplitud sintética con la de un brazo que simula el mecanismo de formación del artefacto. La distancia entre la distribución de amplitudes sintéticas y la medida en las observaciones reales de artefacto es el contraste de realismo. Repite para la apariencia la lógica del Objetivo~2, que compara distribuciones contra una referencia externa y no casos emparejados. La superioridad frente a la inserción por copia y pegado se declara control de cordura y no evidencia de calidad, porque la amplitud de esa línea base vale cero por construcción y cualquier brazo que genere algo la supera. Es una debilidad reconocida del diseño y no una victoria del método. Ningún resultado de ese contraste cuenta a favor del sintetizador, y la decisión del Objetivo~3 recae en la comparación de equivalencia.

Se comparan tres brazos sobre la misma anatomía de origen y las mismas poses de implante: el sintetizador propuesto, la inserción por copia y pegado y el protocolo físico de Peters et al. La inserción por copia y pegado es la línea base ingenua, que coloca metal sin artefacto \GAPDEC{definición operativa de la inserción por copia y pegado, en particular el valor de HU asignado a $M$}. El protocolo físico, implementado con CatSim dentro de XCIST \cite{deman2007catsim,wu2022xcist} y asociado al desafío AAPM CT-MAR \cite{haneda2025aapm}, es el brazo que aporta el mecanismo, porque es el único de los tres que simula la física de formación del artefacto. Se adopta en lugar de una reimplementación validada propia de XCIST/CatSim, que queda fuera de alcance.

El protocolo físico publicado es bidimensional y evalúa MAR, de modo que extenderlo a la síntesis de osteosíntesis pélvica exige una adaptación explícita que no hereda su validación. La geometría con que generan sus datos simula una sola fila de detector con 900 columnas y 1\,000 vistas, con la dispersión calculada para una fila y escalada a 64, y reconstruye por retroproyección filtrada con corrección de agua \cite{peters2025hybrid}. El implante de este trabajo recorre el corredor entero, cuya longitud mediana es de 138~mm (Tabla~\ref{tab:preinscripcion}), así que la adaptación cubre toda la extensión axial del corredor y no un corte aislado. El metal de sus datos de entrenamiento son formas fractales aleatorias insertadas en posiciones aleatorias de tejido blando o hueso, y no implantes \cite{peters2025hybrid}. Peters et al. validan su simulación contra un fantoma físico, sin cubrir el paso híbrido con imágenes clínicas ni esas geometrías genéricas.

Reproducir ese protocolo no es validar el simulador. Validarlo exigiría comparar la simulación contra mediciones físicas en un fantoma con material metálico conocido, escaneado en un tomógrafo real, y esos recursos no están disponibles. Lo que se hizo es una verificación de reproducción: se corrió el protocolo en local y se revisó su código publicado. En ese código, paciente y metal se proyectan juntos, en una sola simulación, sobre un volumen de tres materiales. Son agua con el paciente, agua con la máscara del metal en signo negativo, que desplaza el agua, y la aleación con esa misma máscara. La inserción ocurre, por tanto, por desplazamiento de material antes de la proyección y no por pegado en el dominio de imagen, y eso es lo que hace híbrido al protocolo.

La revisión del código se hizo sobre dos repositorios publicados, en las revisiones \texttt{4cf3544} del simulador y \texttt{4993e87} de los ejemplos, cuya licencia declarada es la de tres cláusulas (BSD 3-Clause). El script de simulación publicado no corre tal como viene: aborta al serializar su propia configuración, de modo que reproducirlo exige un parche de una línea. Ese parche se aplicó sobre una copia local y no sobre el clon del repositorio original, para que el clon revisado quede intacto \GAPLIT{referencia bibliográfica del repositorio publicado del protocolo físico, con las revisiones y la licencia que se revisaron; el código no tiene ficha ni entrada de bibliografía}. El protocolo físico se corre sobre un subconjunto reducido de pacientes \GAPDEC{número de pacientes del subconjunto del protocolo físico}, con un control de simulación y reconstrucción sin metal que aísla los cambios que introduce el propio protocolo \GAPDATO{resultados del protocolo físico de Peters et al. sobre la anatomía y las poses de este trabajo}.

La superioridad sobre la inserción por copia y pegado se evalúa con una prueba de rangos con signo de Wilcoxon de una cola, con nivel de significación de 0.05. La prueba se parea por paciente, porque los dos brazos se aplican a los mismos pacientes. La unidad de agregación es el paciente: la mediana de la amplitud sobre sus regiones de medición y sus cortes deja un valor único. Las pruebas pareadas se calculan sobre los 14 valores de los pacientes de prueba sin metal (Sección~\ref{sec:datos}). Se reportan el tamaño de efecto y su IC además del valor $p$.

Para la comparabilidad con el protocolo físico se usan dos pruebas unilaterales de equivalencia (TOST, \emph{two one-sided tests}). Se concluye equivalencia solo si el IC del 90\,\% de la diferencia pareada cae entero dentro de un margen $[-\Delta, +\Delta]$ fijado de antemano. Nunca se infiere equivalencia de una diferencia no significativa \GAPDEC{si el brazo de equivalencia frente al protocolo físico se mantiene o pasa a trabajo futuro, contingencia de plazo registrada por la autora}. El margen $\Delta$ se define como la variabilidad propia del método bajo condiciones que no deberían cambiar el resultado. Esa variabilidad suma la que producen distintas semillas del sintetizador sobre un mismo caso y la de corridas repetidas del protocolo físico sobre ese caso \GAPDEC{qué cambia entre dos corridas repetidas del protocolo físico}. Ambas se miden solo en los casos de validación de la partición (Sección~\ref{sec:datos}), de modo que el protocolo físico también se corre sobre ellos; tras aplicar el criterio de inclusión, quedan 3 pacientes con implante real. El número de semillas por caso no está fijado (Sección~\ref{sec:sintetizador}).

Fuera de la banda, el RMSE y el índice de similitud estructural (SSIM, \emph{structural similarity index}) contra la TC original comprueban la preservación, restringidos al complemento de $B_{\delta}$. En el sintetizador esa preservación es exacta por construcción (Sección~\ref{sec:sintetizador}); en el protocolo físico sí es una medición, porque la simulación y la reconstrucción pueden alterar la imagen entera. Dentro de $B_{\delta}$, la desviación es la señal buscada y no se puntúa como error.

\subsection{Resumen del diseño experimental}\label{sec:resumen-diseno}

La Tabla~\ref{tab:diseno} liga cada objetivo evaluado con su variable independiente, su variable dependiente, su comparación y su análisis. Los tres fijan de antemano la magnitud que se mide, pero no con el mismo grado de cierre. En el Objetivo~1, el criterio y la regla de decisión se fijaron antes de correr la prueba. En el Objetivo~2 se preinscribieron la distribución de poses y la métrica, pero no una regla de decisión sobre la distancia. En el Objetivo~3 se fijaron el criterio primario, su definición operativa y el orden de los contrastes, y el margen $\Delta$ queda pendiente de preinscripción.

\begin{table}[htbp]
\centering
\small
\caption{Diseño experimental por objetivo.}
\label{tab:diseno}
\begin{tabular}{p{0.12\linewidth} p{0.16\linewidth} p{0.19\linewidth} p{0.17\linewidth} p{0.2\linewidth}}
\hline
\textbf{Objetivo} & \textbf{Variable independiente} & \textbf{Variable dependiente} & \textbf{Comparación} & \textbf{Análisis} \\
\hline
1. Representación & Variante del autoencoder (2) $\times$ codificación (3) & MAE en hueso por paciente, en HU; MAE en $B_{\delta}$, descriptivo & Criterio de 25~HU & Media sobre 34 pacientes de prueba; IC 95\,\% por remuestreo de pacientes; orden a priori \\
2. Colocación & Poses del muestreador preinscrito; cohorte primaria (72) y de sensibilidad (49) & Distribución de grados de brecha (0 a 3); fracción por zona de densidad y viabilidad del corredor por caso, descriptivas & Series navegada y convencional de Zwingmann et al., en S1 & Distancia de Wasserstein-1 en grados, sin regla de decisión; estratificación post hoc por viabilidad del calibre \\
3. Apariencia & Brazo: sintetizador, copia y pegado, protocolo físico & \emph{Streak amplitude} (primaria); \emph{bone} y \emph{metal integrity} (descriptivas) & Protocolo físico; observaciones reales de artefacto; copia y pegado & TOST con margen $\Delta$ de validación, contraste primario; distancia entre distribuciones; Wilcoxon pareado de una cola, control de cordura \\
\hline
\end{tabular}
\end{table}

\section{Amenazas a la validez}\label{sec:amenazas}

Las amenazas se agrupan en cuatro tipos: interna, de constructo, externa y de la conclusión. Para cada una se indica qué se hizo para controlarla o por qué queda abierta.

\textbf{Validez interna.} En el Objetivo~2, la amenaza principal es la circularidad entre muestreador y referencia clínica, y se controla con la preinscripción: ningún parámetro sale de la referencia clínica y la distribución se congeló antes de calcular cualquier distancia contra ella. Tres decisiones, en cambio, se tomaron después de ver datos: la fracción de limpieza de componentes, la regla de ocupación del implante y la estratificación por viabilidad del calibre. El control de nivel excluyó a 11 de 34 pacientes con objetos no ortopédicos y a 7 de 57 sin objeto (Sección~\ref{sec:datos}). La cohorte primaria pierde así, en proporción, más pacientes del grupo~2 que del grupo~3, y su composición por grupos deja de ser la de los 91 pacientes con S1 localizado. Si el corredor difiriera entre grupos, la distribución de grados heredaría esa composición; la cohorte de sensibilidad, que solo reúne pacientes sin objeto, no está expuesta a ese desequilibrio.

La fractura es el confundidor de esa estratificación, y el cribado de la Sección~\ref{sec:datos} lo midió con grupo de comparación. La fractura confirmada aparece por igual en los dos estratos, 10 de 15 corredores estrechos y 10 de 15 controles, así que no explica el estrechamiento. La fractura sacra desplazada, que es la única con mecanismo conocido, sí se concentra en los corredores estrechos: 7 de 15 frente a 2 de 15, con $p = 0.109$. Con 15 casos por grupo esa diferencia es sugerente y no concluyente. El estrechamiento no puede atribuirse entonces ni a anatomía normal ni a patología, y esta muestra no permite cuantificar la fracción que corresponde a cada una.

En el Objetivo~1, la regla operativa y sus seis combinaciones se fijaron después de una exploración del autoencoder preentrenado que incluía a los pacientes de prueba y que ya había fallado el criterio de 25~HU. Ese umbral era anterior a la exploración y no se movió. El rediseño del Objetivo~3 también es posterior al veredicto del Objetivo~1, aunque la regla de la compuerta no se tocó. Lo mismo ocurre con la extensión de la compuerta al latente de Guo et al., cuya cota por saturación se calculó sobre los mismos pacientes de prueba.

\textbf{Validez de constructo.} SAP mide protrusión fuera de una envolvente ósea segmentada, no perforación de una cortical segmentada; su resolución es la del control sobre el fantoma (Sección~\ref{sec:sap}). El cierre morfológico de la envolvente (Sección~\ref{sec:corredor}) podría ocupar el canal sacro o los forámenes, y en ese caso una perforación hacia ellos daría brecha nula \GAPDATO{comprobación de que la envolvente ósea deja fuera el canal sacro y los forámenes sacros}. El marco de referencia se define perpendicular al platillo superior de S1, mientras que el control de nivel y las medidas sobre máscaras usan el techo de la etiqueta de S1 (Sección~\ref{sec:marco-metal}). Ese techo puede estar en la cresta sacra media, así que las dos superficies pueden no coincidir, y esa diferencia no está medida \GAPDATO{diferencia en milímetros entre el techo de la etiqueta de S1, que puede estar en la cresta sacra media, y el platillo superior sobre el que se define el marco de referencia}. Los límites de grado son una convención: el ancho de 2~mm procede de la literatura de tornillos pediculares, de la que Smith et al.~\cite{smith2006iliosacral} declaran tomar la escala, y su traslado al corredor iliosacro es un supuesto explícito.

La referencia clínica no dice a qué grado pertenece una perforación de exactamente 2 o 4~mm, ni reporta el grosor de corte. Tejwani et al.~\cite{tejwani2014} describen más penetración foraminal en cortes de 2.5~mm que en cortes de 5.0~mm, aunque con $p = 0.3$. Cualquier acuerdo con la referencia clínica se lee como acuerdo entre distribuciones de grados medidas a una resolución de referencia desconocida, no como equivalencia vóxel a vóxel. El nivel de la referencia es un supuesto en una de sus dos series. Zwingmann et al. nombran S1 en la técnica de la serie convencional, pero no en la de la navegada. La definición de posición ideal del estudio se refiere al platillo sacro correspondiente y a los forámenes de S1. Este trabajo asume S1 también para la serie navegada, y la publicación no permite verificarlo.

La Ecuación~\eqref{eq:w1} trata los grados como equiespaciados, lo que es también una convención, y la convención $h = 2\sigma$ fija la escala de todas las perturbaciones y, con ella, la distancia resultante. En la apariencia, las métricas de Peters et al. se diseñaron para reducir artefacto y no para generarlo, y una máscara binaria no codifica la aleación del implante. La banda $B_{\delta}$ trunca a propósito las rayas lejanas (Sección~\ref{sec:sintetizador}), y el protocolo físico, que simula y reconstruye la imagen entera, no las trunca. La comparación entre ambos no puede mostrar, por tanto, si el sintetizador reproduciría las rayas más allá de la banda. El sintetizador se entrena, además, con componentes metálicos de todo tipo bajo el supuesto no verificado de que la relación entre máscara y artefacto no depende del tipo de implante (Sección~\ref{sec:sintetizador}). Por último, el trabajo asume que la apariencia de un artefacto metálico puede generarse en el dominio de imagen sin datos de proyección. Lin et al.~\cite{lin2019} y Li et al.~\cite{li2024} reportan una penalización por operar solo en imagen, pero en reducción de artefactos. Lin et al. describen esa reducción como un problema mal planteado, porque paciente y metal contribuyen a la misma traza en el sinograma; la síntesis no tiene que separar esas dos contribuciones. La comparación con el protocolo físico es la que puede contradecir ese supuesto.

El suelo de $-1000$~HU de la representación (Sección~\ref{sec:sintetizador}) censura la cola inferior del campo de desviación, y el sesgo que introduce es siempre a la baja. Ese sesgo no pesa igual en los dos contrastes que miden la amplitud. Frente al protocolo físico apenas pesa, porque su salida reconstruida baja poco por debajo de ese suelo. Frente a la distribución de las observaciones reales de artefacto pesa mucho más, porque esos volúmenes alcanzan valores de TC muy inferiores. Esa diferencia se apoya en una sola corrida local del protocolo físico, de modo que ordena la gravedad del sesgo y no fija una cota. La decisión que trata el suelo compara la pérdida atribuible a la representación con el margen $\Delta$ y lo declara limitación de alcance solo si esa pérdida es menor \GAPDEC{contra cuál de los dos contrastes, el protocolo físico o la distribución de las observaciones reales, se evalúa esa regla}.

\textbf{Validez externa.} La referencia clínica proviene de pelvis fracturadas (tipos~B y~C de Tile), que es la población en la que el procedimiento se practica, mientras que los corredores de este trabajo se miden en pelvis sin osteosíntesis. Reilly et al.~\cite{reilly2003effect} midieron en seis pelvis cadavéricas que un desplazamiento craneal de 5 a 20~mm de una fractura sacra de zona~II, creada por osteotomía, reduce entre 36 y 90\,\% el área disponible para tornillos iliosacros en S1. Una malreducción residual estrecharía, por ese mecanismo, los corredores de la serie clínica. La anatomía receptora tampoco está libre de fractura: el cribado de fractura halló fractura confirmada en 20 de los 30 casos que examinó (Sección~\ref{sec:datos}). Eso acerca la anatomía receptora a la población en que el procedimiento existe, en vez de alejarla. Lo que no puede afirmarse es que las dos poblaciones estén emparejadas, porque la fractura de la cohorte receptora es incidental, no está graduada ni clasificada por un especialista y su estado de reducción es desconocido. Además, la cohorte local es una fracción de la colección publicada, su reconstrucción es desconocida, la publicación de TotalSegmentator no reporta exactitud para la etiqueta de S1, y el protocolo físico adoptado se validó en dos dimensiones y para MAR. En el Objetivo~3, los tres desplazamientos de dominio entre entrenamiento y síntesis (Sección~\ref{sec:sintetizador}) limitan también la transferencia de lo aprendido sobre metal real a los tornillos paramétricos.

\textbf{Validez de la conclusión.} Las dos series de la referencia clínica son pequeñas (26 y 35 tornillos, Sección~\ref{sec:sap}) y cuentan tornillos, no pacientes independientes. La distancia de Wasserstein-1 se reporta sin prueba inferencial, porque la preinscripción no prevé ninguna, y su estabilidad solo se aprecia al comparar las dos cohortes \GAPDEC{si se añade a la distancia de Wasserstein-1 un intervalo por remuestreo de pacientes, como análisis no preinscrito}. En la compuerta del Objetivo~1, la multiplicidad no corregida solo puede favorecer un veredicto aprobatorio, porque basta con que apruebe una de seis combinaciones; como el veredicto fue negativo, esa multiplicidad no lo compromete. La extensión posterior al latente de Guo et al. añade un séptimo candidato, examinado sobre los mismos pacientes de prueba. En la apariencia, el margen $\Delta$ descansa en solo 3 pacientes de validación (Sección~\ref{sec:apariencia}); las semillas multiplican las observaciones por paciente, pero no el número de pacientes. La prueba de equivalencia tiene además muestras pequeñas, acotadas por el subconjunto del protocolo físico (Sección~\ref{sec:apariencia}), de modo que puede no alcanzar a concluir equivalencia ni diferencia.
